\section{Experiments}
\label{sec:experiments}

% \subsection{Setup}

\mysection{Models and Domains.} We evaluate on Qwen2-7B-Instruct~\cite{yang2024qwen2} and LLaMA-3-8B-Instruct~\cite{dubey2024llama} across 22 datasets spanning QA, commonsense reasoning, math, knowledge, classification, code, and instruction-following (Appendix~\ref{app:domains}). Each dataset is fine-tuned with LoRA (rank 16, $\alpha{=}32$) and merged.

\paragraph{Safety Benchmarks.}
We evaluate on three harmful-prompt benchmarks: \textbf{BeaverTails}~\cite{ji2024beavertails} (700 prompts, 14 categories) is used as the harmful data source during \ssn{} training (\S\ref{sec:gt}); the evaluation split is disjoint from training. \textbf{AdvBench}~\cite{zou2023universal} (520 prompts) and \textbf{HEx-PHI}~\cite{qi2024finetuning} (330 prompts, 11 categories) are never seen in any training phase, serving as zero-shot safety tests.

\paragraph{Evaluation Protocol.}
We use three holdout splits (Table~\ref{tab:splits}), each removing four datasets from different task categories. The hypernetwork trains on the remaining 18 datasets and is evaluated on held-out datasets it has never seen. We report \textbf{Task Accuracy} (domain-specific test performance; official HuggingFace validation/test splits where available, 20\% holdout otherwise) and \textbf{Harmful Rate} (percentage of harmful prompts producing harmful output, lower is safer; 700 BeaverTails prompts disjoint from training, judged by the beaver-dam-7b moderation model~\cite{ji2024beavertails}). We compare against the \textbf{base model} (aligned, no fine-tuning) and the \textbf{LoRA model} (fine-tuned, no safety mechanism). Table~\ref{tab:comparison} provides a qualitative comparison with prior defense methods.

\begin{table}[h]
\caption{Holdout splits. Each split removes four datasets from different task categories. $^\dagger$20\% holdout eval (no standard test set).}
\label{tab:splits}
\centering
\renewcommand{\arraystretch}{1.4}
\resizebox{\columnwidth}{!}{%
\begin{tabular}{@{}cllll@{}}
\toprule
\textbf{Split} & \textbf{Domain 1} & \textbf{Domain 2} & \textbf{Domain 3} & \textbf{Domain 4} \\
\midrule
A & BoolQ {\small(QA)} & ARC {\small(Science)} & Comp.\ Math {\small(Math)} & Alpaca$^\dagger$ {\small(Instr.)} \\[3pt]
B & PIQA {\small(Physical)} & HellaSwag {\small(Common.)} & CSenseQA {\small(Common.)} & OpenHermes$^\dagger$ {\small(Instr.)} \\[3pt]
C & MMLU {\small(Knowledge)} & NQ-Open {\small(QA)} & AG News$^\dagger$ {\small(Classif.)} & Dolly$^\dagger$ {\small(Instr.)} \\
\bottomrule
\end{tabular}%
}
\end{table}

\subsection{Main Results}

Table~\ref{tab:main_results} shows results across all three splits. LoRA fine-tuning increases the harmful rate from 3--7\% (base) to as high as 43\%, with the largest degradation on instruction-following and commonsense domains. \method{} reduces the harmful rate to below 1\% across all 12 holdout datasets and both architectures, while preserving task accuracy within 1 percentage point. The highest residual rates (NQ-Open, 0.6\%) still represent $>$95\% relative reduction from the fine-tuned model.

% Combined main results — LLaMA-3-8B + Qwen2-7B, LoRA vs OneShot vs HyperSafe.
\begin{table}[t]
\centering
\renewcommand{\arraystretch}{1.1}
\resizebox{\textwidth}{!}{%
\begin{tabular}{@{}l ccc ccc ccc ccc@{}}
\toprule
& \multicolumn{3}{c}{\textbf{Task Acc (\%) $\uparrow$}}
& \multicolumn{3}{c}{\textbf{LoRA Harm (BT/AB/HX, \%) $\downarrow$}}
& \multicolumn{3}{c}{\textbf{OneShot Harm $\downarrow$}}
& \multicolumn{3}{c}{\textbf{\method{} Harm $\downarrow$}} \\
\cmidrule(lr){2-4} \cmidrule(lr){5-7} \cmidrule(lr){8-10} \cmidrule(lr){11-13}
\textbf{Domain} & LoRA & OneShot & \method{} & BT & AB & HX & BT & AB & HX & BT & AB & HX \\
\midrule
\multicolumn{13}{@{}l}{\textit{\textbf{LLaMA-3-8B}}} \\
\midrule
BoolQ                  & 87.2 & 87.2 & 86.9 &  9.7 &  1.5 &  7.0 & 0.3 & 0.0 & 1.0 & 0.1 & 0.0 & 0.0 \\
ARC                    & 86.6 & 18.8 & 86.5 & 41.7 & 33.8 & 20.0 & 0.1 & 0.0 & 0.7 & 0.0 & 0.0 & 0.0 \\
Comp.\ Math$^\dagger$  & 33.8 & 28.4 & 34.1 &  8.6 &  1.7 &  6.0 & 0.9 & 0.4 & 1.7 & 0.0 & 0.0 & 0.0 \\
Alpaca$^\dagger$       & 39.0 & 15.2 & 38.9 & 23.0 &  9.6 & 34.3 & 0.1 & 0.0 & 0.7 & 0.1 & 0.0 & 0.0 \\
PIQA                   & 88.1 & 87.2 & 87.7 & 37.7 & 40.4 & 37.7 & 0.4 & 0.0 & 0.7 & 0.0 & 0.0 & 0.0 \\
HellaSwag              & 89.4 & 87.0 & 89.1 & 12.1 & 15.2 & 11.3 & 0.4 & 0.4 & 0.7 & 0.1 & 0.0 & 0.3 \\
CSenseQA               & 81.7 & 79.9 & 81.5 & 17.7 &  5.4 &  5.0 & 0.0 & 1.9 & 6.7 & 0.0 & 0.0 & 0.0 \\
OpenHermes$^\dagger$   & 54.9 & 16.4 & 54.6 & 41.0 & 85.6 & 62.0 & 0.0 & 0.0 & 2.7 & 0.0 & 0.0 & 0.0 \\
MMLU                   & 61.1 & 61.5 & 61.0 &  9.9 & 10.2 &  8.7 & 0.3 & 0.0 & 1.0 & 0.0 & 0.0 & 0.0 \\
NQ-Open$^\dagger$      & 77.4 & 55.9 & 77.1 & 20.1 & 11.5 & 19.7 & 2.7 & 0.2 & 2.0 & 0.4 & 0.0 & 0.0 \\
AG News$^\dagger$      & 94.1 & 93.9 & 90.6 &  9.0 &  0.8 &  5.3 & 2.7 & 2.3 & 9.3 & 0.0 & 0.0 & 0.0 \\
Dolly$^\dagger$        & 37.8 & 16.9 & 34.4 & 38.6 & 64.0 & 49.0 & 0.1 & 0.0 & 0.7 & 0.1 & 0.0 & 0.0 \\
\textit{Avg.}          & \textit{69.3} & \textit{54.0} & \textit{68.5}
                       & \textit{22.4} & \textit{23.3} & \textit{22.2}
                       & \textit{0.7}  & \textit{0.4}  & \textit{2.3}
                       & \textit{0.1}  & \textit{0.0}  & \textit{0.0} \\
\midrule
\multicolumn{13}{@{}l}{\textit{\textbf{Qwen2-7B}}} \\
\midrule
BoolQ                  & 88.4 & 88.5 & 88.2 & 14.3 &  7.5 & 13.3 &  4.4 & 0.2 &  2.3 & 0.1 & 0.0 & 0.0 \\
ARC                    & 91.3 & 91.5 & 91.3 &  7.1 &  0.2 & 33.7 &  1.0 & 0.0 &  3.7 & 0.0 & 0.0 & 0.0 \\
Comp.\ Math$^\dagger$  & 74.4 & 76.1 & 74.1 &  6.0 &  0.0 &  6.7 &  5.4 & 0.0 &  8.3 & 0.0 & 0.0 & 0.0 \\
Alpaca$^\dagger$       & 39.6 & 37.2 & 39.5 & 19.4 & 13.3 & 33.0 &  5.4 & 1.0 &  8.0 & 0.4 & 0.0 & 0.0 \\
PIQA                   & 90.8 & 90.6 & 90.6 & 43.1 & 67.7 & 57.3 &  4.3 & 0.2 &  2.0 & 0.0 & 0.0 & 0.0 \\
HellaSwag              & 93.2 & 58.0 & 92.8 & 32.9 & 29.8 & 45.3 &  4.7 & 1.0 &  5.0 & 0.1 & 0.0 & 0.3 \\
CSenseQA               & 84.6 & 84.2 & 84.5 & 15.7 &  3.7 & 14.0 &  1.0 & 0.0 &  2.0 & 0.0 & 0.0 & 0.0 \\
OpenHermes$^\dagger$   & 57.2 & 54.6 & 57.0 & 38.0 & 70.2 & 64.7 &  1.1 & 3.3 &  8.3 & 0.1 & 0.0 & 0.1 \\
MMLU                   & 68.7 & 68.7 & 68.4 &  9.7 &  0.2 & 22.0 &  5.1 & 0.2 &  6.3 & 0.1 & 0.0 & 0.3 \\
NQ-Open$^\dagger$      & 54.8 & 23.7 & 54.3 & 11.9 &  6.0 & 28.7 & 11.9 & 1.0 & 14.7 & 0.6 & 0.0 & 0.7 \\
AG News$^\dagger$      & 97.6 & 92.5 & 97.6 & 10.1 &  0.2 &  6.0 &  3.7 & 0.0 &  6.7 & 0.0 & 0.0 & 0.0 \\
Dolly$^\dagger$        & 35.6 & 33.8 & 35.4 & 22.1 & 24.0 & 44.3 & 11.0 & 2.7 & 11.7 & 0.3 & 0.0 & 0.0 \\
\textit{Avg.}          & \textit{73.0} & \textit{66.6} & \textit{72.8}
                       & \textit{19.2} & \textit{18.6} & \textit{30.8}
                       & \textit{4.9}  & \textit{0.8}  & \textit{6.6}
                       & \textit{0.1}  & \textit{0.0}  & \textit{0.1} \\
\bottomrule
\end{tabular}%
}
\caption{Main results across 12 holdout domains on LLaMA-3-8B and Qwen2-7B. \method{} reduces harmful rates to near-zero on all three benchmarks (BeaverTails, AdvBench, HEx-PHI) while preserving accuracy within ${<}1$~percentage~point of the LoRA fine-tuned baseline. OneShot~\citep{kim2024oneshot} achieves comparable safety on most LLaMA-3 domains but catastrophically degrades accuracy on instruction-following datasets (ARC, Alpaca, OpenHermes, Dolly); on Qwen2 it leaves substantial residual harm (avg 6.6\% HEx-PHI). $^\dagger$20\% holdout eval (no standard test set).}
\label{tab:main_results}
\end{table}


\paragraph{Comparison with Llama Guard 4.}
We compare against Llama Guard 4 (12B), the latest publicly released version of the Llama Guard family (originally introduced by~\citet{inan2023llama}; Llama Guard 4 released as part of the Llama 4 suite~\citep{meta2025llama4}), on the same three benchmarks (Appendix Table~\ref{tab:llamaguard}). On the two fully zero-shot benchmarks --- the appropriate comparison, since BeaverTails appears in \ssn{} training (with a disjoint eval split) --- Llama Guard 4 misses 6.5\% of AdvBench and 6.7\% of HEx-PHI prompts, while \method{} misses 0.0\% and 0.3\% respectively. This is achieved with an \ssn{} of only $\sim$3\% of backbone parameters versus Llama Guard 4's dedicated 12B model. The gap is wider on BeaverTails (53.6\% vs.\ 0.7\%).

\paragraph{Ablations and additional analyses.}
Five ablations are reported in full in Appendices~\ref{app:ablation_summary}--\ref{app:additional}: (i)~loss components, (ii)~training-data diversity (instruction-following is the most important category), (iii)~calibration set size (saturates at $N_\text{cal}{=}50$), (iv)~generalization to full fine-tuning (\ssn{}s trained for LoRA models transfer), and (v)~cross-benchmark safety transfer ($<$1\% harm on the zero-shot AdvBench and HEx-PHI benchmarks).